\documentclass[review,onefignum,onetabnum]{siamart251216}

% Mathematical symbols. The SIAM class already loads amsmath.
\usepackage{amssymb}

% Tables and figures.
\usepackage{booktabs}
\usepackage{graphicx}

% Oxford comma in multiple cleveref references.
\newcommand{\creflastconjunction}{, and~}

% The theorem, lemma, corollary, proposition, definition, and proof
% environments are already defined by siamart251216.
\newsiamremark{remark}{Remark}
\newsiamremark{observation}{Observation}

% Custom mathematical notation.
\newcommand{\fl}{\operatorname{fl}}
\newcommand{\R}{\mathbb{R}}
\newcommand{\norm}[1]{\left\lVert #1 \right\rVert}

\emergencystretch=1em

\newcommand{\submissiondate}{August 03, 2026}

\title{Backward Error of the Matrix Inverse in the
Purcell--Egerv\'ary--Castillo Pivoting Transformation%
\thanks{Submitted to the editors \submissiondate.}}

% Running headers. The abbreviated title must not exceed 50 characters.
\headers{Backward Error of a Matrix Inverse}{F. R. Villatoro}

\author{Francisco R. Villatoro%
\thanks{Escuela de Ingenier\'ias Industriales, Universidad de M\'alaga,
29071 M\'alaga, Spain
(\email{frvillatoro@uma.es}). \funding{This work was  supported by the Universidad de M\'alaga, Spain,
through project PPRO-B4-2026-004 of its Plan Propio de Investigaci\'on, Transferencia y Divulgaci\'on Cient\'ifica.}}}

\begin{document}

\maketitle

\begin{abstract}
The pivoting transformation of Castillo, Cobo, Jubete, Pruneda, and Castillo
[SIAM J. Matrix Anal. Appl., 22 (2000), pp.~666--681] computes a matrix inverse
through a sequence of interpretable tableaux. Algebraically the transformation is
the vector method of Purcell (1953), rediscovered repeatedly as an implicit LU or
direct projection method; but the backward error of the explicitly formed inverse
has not been analyzed. We prove the exact residual decomposition
$R_N=D_N-AE_{N+1}$ into a transformation defect, created when the elementary
transformations are formed from rounded quantities, and an accumulated application
error of Wilkinson type. Thus a bound on the latter alone is not a backward-error
bound for the inverse. We further derive a row-wise representation of $D_N$, a
growth factor with an unavoidable linear floor in the number of steps, and an exact
growth bound for strictly row-diagonally dominant matrices. Two numerical campaigns,
including $14.7$ million matrices, show no observed violation of the rigorous growth-based estimate on the
reliability-screened comparisons, while it overestimates by a median
factor $2.3\times10^{5}$, with the pessimism overwhelmingly concentrated in the
accumulation term. The transformation defect is the sharpest ex-post diagnostic
measured, while the minimum scaled pivot is the dominant local indicator of its
formation. A paired replay shows that row pivoting more than doubles the defect on
two diagonally dominant ensembles while changing the application term only
marginally; the effect on the final residual is strongly modulated by the baseline
weight of the defect. Finally, on an orthogonal family with $\kappa_2(A)=1$ the
first-nonzero rule produces tableau growth asymptotic to $\sqrt2/\varepsilon$, and
the campaign records inverse backward error exceeding $10^7u$ at unit condition
number.
\end{abstract}

\begin{keywords}
matrix inversion, backward error analysis, floating-point arithmetic,
pivoting, biorthogonalization, numerical stability
\end{keywords}

\begin{MSCcodes}
65F05, 65G50, 65F35, 65F25
\end{MSCcodes}

\section{Introduction}
\label{sec:intro}

Let $A \in \R^{n\times n}$ be nonsingular. The pivoting transformation introduced by
Castillo, Cobo, Jubete, Pruneda, and Castillo \cite{CastilloCJPC2000,CastilloCJP1999}
constructs a sequence of tableaux $V^1,\dots,V^{n+1}$ by the elementary rule
%
\begin{equation}\label{eq:castillo-pivot}
v^{j+1}_k =
\begin{cases}
v^j_k / t^j_j, & k = j,\\[2pt]
v^j_k - \dfrac{t^j_k}{t^j_j}\, v^j_j, & k \neq j,
\end{cases}
\qquad
(t^j)^T = a_j^T V^j ,
\end{equation}
%
where $a_j^T$ denotes the $j$-th row of $A$. Starting from $V^1 = I$ and using the rows of $A$
in order, the process terminates with $V^{n+1} = A^{-1}$. Its distinctive feature is the interpretation of the intermediate tableaux: each $V^j$ contains
the generators of the subspace orthogonal to $\{a_1,\dots,a_{j-1}\}$ together with those of its
complement, so the inverses and determinants of all leading principal submatrices come without
extra computation, and the inverse can be updated after a change of one row by a single further
pivoting step.

The algebraic theory of~\eqref{eq:castillo-pivot} was developed in~\cite{CastilloCJPC2000} from
the point of view of orthogonal decompositions and applied there to matrix inversion, determinant
evaluation, rank determination, linear dependence testing, subspace intersection, and the
compatibility analysis of linear systems. The numerical behavior of the process was explicitly
left open: in the words of the authors~\cite[p.~667]{CastilloCJPC2000}, its numerical properties
with respect to stability and ill conditioning \emph{``will be the aim of another paper''}. To the
best of our knowledge that paper never appeared. Six years later the same group returned to the
question only qualitatively, warning against the assumption that orthogonalization procedures are
safe and recommending standard pivoting and preconditioning
techniques~\cite[\S2.1]{CastilloPSM2006,HerreroCP2002}. No backward error bound was derived in any
of these works, and the present paper supplies part of the missing analysis. Before stating what
we do, we place~\eqref{eq:castillo-pivot} in a literature which is unusually scattered, and in
which the same algorithm has been discovered independently at least a dozen times.

\subsection{Genealogy and the stability question}
\label{sec:genealogy}

Algebraically,~\eqref{eq:castillo-pivot} is Purcell's \emph{vector method}~\cite{Purcell1953},
identified by Egerv\'ary as a case of rank reduction~\cite{Egervary1957,Egervary1960},
developed independently by Hestenes~\cite{Hestenes1958} as inversion by
biorthogonalization, incorporated into the ABS class~\cite{AbaffyBS1984,AbaffySpedicato1989},
and rediscovered in the sparse-solver literature as a direct projection method
\cite{BenziMeyer1995,BenziTuma1998}.  SM1 records the extended genealogy.  Two historical
observations are central here.  First, Hestenes noted that the choice $V^{(0)}=I$ is
numerically fragile because the pivots are tied to the leading principal minors of $A$;
\S\ref{sec:tableaux} makes this precise by identifying the intermediate tableaux with
bordered inverses of the leading principal blocks.  Second, the ABS literature developed
sensitivity and forward-error analyses~\cite{Broyden1985,Galantai1991,AsadbeigiPBJ2017}
but did not obtain a backward-error bound for the explicitly formed inverse.

The closest backward analysis is Fletcher's~\cite{Fletcher1997} treatment of the LIU
factorization under partial pivoting.  It controls a perturbation of the factorization but
never forms the inverse, and Fletcher explicitly notes the difficulty of transferring the
partial-sum errors to the data in a Wilkinson-style solution analysis.  In the opposite
regime, G\'ati~\cite{Gati2012,Gati2013}, extending the automatic roundoff analysis of
Miller and Spooner~\cite{MillerSpooner1978a}, found large normwise errors for the
unpivoted implicit-LU method at condition number $12.3$.  The apparent conflict is
therefore a pivoting issue: Fletcher assumes partial pivoting, whereas G\'ati's examples
use none.  The Castillo implementation~\cite[Algorithm~1, Step~3]{CastilloPSM2006}
selects a pivot column but, in its original form, processes rows in order and hence lies
between these regimes.

\subsection{Contribution and outline}
\label{sec:contribution}

The missing object in the literature is the backward error of the \emph{explicitly formed inverse}.
Writing $\widetilde V_{N+1}$ for the exact product of the computed elementary transformations and
$E_{N+1}=\widehat V_{N+1}-\widetilde V_{N+1}$, we prove
\begin{equation}\label{eq:intro-decomposition}
R_N=I-A\widehat V_{N+1}=D_N-AE_{N+1},
\qquad
D_N=I-A\widetilde V_{N+1},
\end{equation}
and, for nonsingular $\widehat V_{N+1}$,
$\Delta A=R_N\widehat V_{N+1}^{-1}$.  The first contribution is therefore to identify and analyze
$D_N$, which is absent from a Wilkinson-style accumulation estimate.  Theorem~\ref{thm:defect-structure}
telescopes it from local formation defects, and Proposition~\ref{prop:local-defect} exposes the
scaled pivot and dot-product sensitivity as its natural local variables.

The second contribution is a growth factor $\Gamma_N$ satisfying
$r_R^\infty\le\gamma_{n+3}\Gamma_N$ and the universal floor
$\Gamma_N\ge N/(1+N\gamma_2)$.  The third is an exact matrix-class result: for strictly row-
diagonally dominant $A$ with gap $\alpha$ and $\|A\|_\infty=1$,
Theorem~\ref{thm:dd} gives $\Gamma_N\le6n/\alpha^3$ and shows that first-nonzero and
maximum-modulus follow the same pivot path; Corollary~\ref{cor:dd-fp} gives the corresponding
first-order floating-point residual bound.  The fourth is an orthogonal adversarial family with
$\kappa_2(A)=1$ on which first-nonzero has growth asymptotic to $\sqrt2/\varepsilon$, while
magnitude-based rules have growth $(1-\varepsilon)^{-1/2}$.

The campaigns show that the rigorous bound is valid but highly non-sharp on the screened records,
that $D_N$ is the tightest ex-post diagnostic measured, and that the minimum scaled pivot is the
dominant local trajectory indicator.  They also show that row pivoting is class dependent: useful
on unstructured prescribed-singular-value matrices but harmful on the diagonal-dominance ensembles.
SM1--SM11 contain the extended genealogy, proof details, campaign provenance, reliability
accounting, regression, per-family results, replay analyses, adversarial tables, and high-precision
validation.


\section{Floating-point analysis}
\label{sec:fp}

Throughout, $\fl(\cdot)$ denotes evaluation in binary floating-point arithmetic with unit
roundoff $u$, and
%
\begin{equation}\label{eq:gamma}
\gamma_k = \frac{ku}{1-ku}, \qquad ku<1,
\end{equation}
%
is the standard constant of~\cite[Ch.~3]{Higham2002}. Absolute values and inequalities
between matrices are understood entrywise. When a matrix norm is left unspecified, it is
submultiplicative, satisfies $\|X\|\le\bigl\|\,|X|\,\bigr\|$, and is monotone on nonnegative matrices:
$0\le C\le D$ implies $\|C\|\le\|D\|$.  These properties are sufficient to pass from the
componentwise bounds below to norm bounds and are satisfied by the spectral, Frobenius, and
infinity norms used in the paper. We write
$\widehat X$ for a computed quantity
and $\widetilde X$ for the result of applying computed data in exact arithmetic. Unless stated
otherwise, $A$ denotes the matrix as supplied to the working arithmetic after any initial data
conversion; an additional perturbation caused by rounding an external real-valued data set into
the working format is not included in the model below. All floating-point inequalities below assume
that no underflow or overflow occurs and that every $\gamma_k$ appearing in the statement is
defined, i.e., $ku<1$; individual results state only additional restrictions when needed.

The analysis separates three sources of error, which the evidence of \S\ref{sec:experiments}
shows behave very differently: the rounding committed when each elementary transformation is
\emph{formed} from inner products, reciprocals, and quotients, which produces the
\emph{transformation defect} $D_N$ of \S\ref{sec:defect}; the rounding committed when each
transformation is \emph{applied} to the current tableau, which accumulates into $E_{N+1}$ and
is the object of the classical Wilkinson-style estimate of \S\ref{sec:accumulation}; and the
amplification of both by $\widehat V_{N+1}^{-1}$ when the residual is converted into a backward
perturbation of $A$.

\subsection{Rounded inner products}
\label{sec:inner}

The standard estimate~\cite[Lem.~3.1, Eq.~(3.4)]{Higham2002} states that for $x,y\in\R^n$
%
\begin{equation}\label{eq:inner}
\fl(x^Ty) = (x+\delta x)^Ty, \qquad |\delta x| \le \gamma_n |x|,
\end{equation}
%
and in particular $|\fl(x^Ty)-x^Ty| \le \gamma_n |x|^T|y|$. Applied to the pivot row
of~\eqref{eq:castillo-pivot} this gives, for the computed dot products
$\hat t^{\,j}_k = \fl(a_j^T\widehat v^{\,j}_k)$,
%
\begin{equation}\label{eq:dots}
\bigl|\hat t^{\,j}_k - a_j^T\widehat v^{\,j}_k\bigr|
\;\le\;
\gamma_n\,|a_j|^T\bigl|\widehat v^{\,j}_k\bigr|,
\qquad k=1,\dots,n .
\end{equation}
%
The backward perturbation $\delta x$ in~\eqref{eq:inner} depends on $y$, so when the $n$ inner
products $\hat t^{\,j}_k$ share the argument $a_j$ their individual backward interpretations
need not provide a single common perturbation of $a_j$ satisfying the same componentwise bounds.
Algebraically, if the current tableau is nonsingular, one can force a common row perturbation by
multiplying the vector of dot-product errors by the inverse tableau, but this may amplify the
perturbation by its conditioning and gives no Wilkinson-size backward bound. This is the
obstruction relevant here and the reason the defect below is defined at the level of the computed
coefficients rather than identified a priori with a small perturbation of the data. Two further
roundings occur after the selected column has been exchanged into the active position: the
reciprocal of the selected computed pivot and the quotients by that pivot, each contributing a
relative error at most $u$. Their combination with~\eqref{eq:dots} produces the formation defect
analyzed in \S\ref{sec:defect-structure}.

\subsection{One computed pivoting transformation}
\label{sec:onestep}

Let $\widehat V_j$ be the computed tableau at the start of step $j$. To cover the row-pivoted
variants of \S\ref{sec:experiments} we let $\pi_j$ denote the index of the row of $A$
processed at step $j$ ($\pi_1,\dots,\pi_N$ distinct, $\pi_j=j$ when rows are taken in order),
and to keep the notation of \S\ref{sec:inner} we write $a_j := a_{\pi_j}$ for the row
processed at step $j$, reserving the subscript for the step. The computed dot products
$\hat t^{\,j} = \fl\bigl(a_j^T\widehat V_j\bigr)$ are formed first; the pivot column
$p_j\ (\ge j)$ is then selected among the admissible indices $j,\dots,n$ as a function of the
\emph{computed} values $\hat t^{\,j}$ (and, for the scaled rules, of the computed column
norms), so that every selection rule of \S\ref{sec:experiments}, and every diagnostic
derived from the pivots, operates on rounded quantities; $P_j$ is the elementary permutation
that exchanges columns $j$ and $p_j$. Set
$\overline V_j=\widehat V_jP_j$, write $\overline v_k^{\,j}=\overline V_je_k$, and let
$\overline t^{\,jT}=\hat t^{\,jT}P_j$ denote the computed dot products in the same permuted
column order. Thus $\overline t_j^{\,j}=\hat t_{p_j}^{\,j}$ is the selected computed pivot. The
coefficients actually used by the step are
$\hat\beta=\fl(1/\overline t_j^{\,j})$ and
$\hat\alpha_k=\fl(\overline t_k^{\,j}/\overline t_j^{\,j})$ for $k\ne j$. Write
%
\begin{equation}\label{eq:Bhat}
\widehat B_j = I + e_j\bigl(\hat r^{\,T}-e_j^T\bigr),
\qquad
\hat r = \bigl(-\hat\alpha_1,\dots,-\hat\alpha_{j-1},\ \hat\beta,\
-\hat\alpha_{j+1},\dots,-\hat\alpha_n\bigr)^T ,
\end{equation}
%
so that $\widehat B_j$ is the identity with its $j$-th row replaced by $\hat r^T$. This is
the computed analogue of $M_j^{-1}$ in~\cite[Eq.~(2.3)]{CastilloCJPC2000}.

In exact arithmetic the step enforces $a_j^TV_{j+1}=e_j^T$, and the relations already achieved
persist, since for $i<l$ the permutation $P_l$ fixes coordinate $i$ (since $p_l\ge l>i$) and row $i$
of $B_l$ is $e_i^T$, so $a_i^TV_{l+1}=a_i^TV_lP_lB_l=e_i^TP_lB_l=e_i^T$. Hence after step $j$
the invariant $a_i^TV_{j+1}=e_i^T$ holds for all $i\le j$, and at termination
$\Pi A V_{N+1}=I$ with $\Pi=\sum_{j}e_je_{\pi_j}^T$, that is,
$V_{N+1}=(\Pi A)^{-1}=A^{-1}\Pi^T$. To make the row-pivot convention unambiguous, set
$\bar A=\Pi A$ and $\bar b=\Pi b$, and let $T$ denote any exact or computed terminal tableau for
$\bar A$. The inverse of the original matrix represented by that tableau is $X=T\Pi$, and
\[
I-AX=\Pi^T(I-\bar A T)\Pi,
\qquad
(I-AX)X^{-1}=\Pi^T(I-\bar A T)T^{-1},
\qquad
Xb=T\bar b.
\]
Thus the spectral, infinity, and Frobenius residual norms and the corresponding normwise backward
errors are identical whether they are evaluated for $(A,b,X)$ or for $(\bar A,\bar b,T)$; also
$\kappa_2(\bar A)=\kappa_2(A)$. In theoretical formulas below we therefore suppress the bars and
write $(A,b,V)$ for the row-permuted problem when row pivoting is used. Implementation-facing
identities occasionally restore $\Pi$ explicitly, for example
$D_N=I-\Pi A\widetilde V_{N+1}$ and $R_N=I-\Pi A\widehat V_{N+1}$.

\begin{lemma}[One step]\label{lem:onestep}
The computed tableau produced by one step of~\eqref{eq:castillo-pivot} satisfies
%
\begin{equation}\label{eq:onestep}
\widehat V_{j+1} = \widehat V_j P_j \widehat B_j + F_j,
\qquad
|F_j| \le \gamma_2\,\bigl|\widehat V_jP_j\bigr|\,\bigl|\widehat B_j\bigr| .
\end{equation}
\end{lemma}

\begin{proof}
For $k\ne j$ the computed column
$\fl(\overline v_k^{\,j}-\hat\alpha_k\overline v_j^{\,j})$ involves one multiplication and one
subtraction, so its entries differ from the exact ones by at most
$\gamma_2(|\overline v_{ik}^{\,j}|+|\hat\alpha_k||\overline v_{ij}^{\,j}|)$, while for $k=j$
the computed column $\fl(\hat\beta\overline v_j^{\,j})$ has entrywise error at most
$u|\hat\beta||\overline v_{ij}^{\,j}|$. Collecting and comparing with the entries of
$|\widehat V_jP_j||\widehat B_j|=|\overline V_j||\widehat B_j|$ gives~\eqref{eq:onestep}.
\end{proof}

\begin{remark}\label{rem:no-backward-step}
It is tempting to rewrite~\eqref{eq:onestep} as
$\widehat V_{j+1} = (\widehat V_jP_j + F_j\widehat B_j^{-1})\widehat B_j$ and read it as a
backward-stable step. We do not. For
$\widehat B_j=I+e_j(\hat r^T-e_j^T)$, the nontrivial row of $\widehat B_j^{-1}$ has entries
$\hat\beta^{-1}$ at the pivot and $\hat\alpha_k/\hat\beta$ off the pivot, quantities of the
scale of the computed dot products rather than $|\overline t_j^{\,j}|^{-1}$. The forward error
in~\eqref{eq:onestep}, however, is proportional to $|\widehat B_j|$; after the backward
reinterpretation it is therefore controlled by a product of the form
$\|\widehat B_j\|\,\|\widehat B_j^{-1}\|$, which can be large when the active pivot is small
relative to the other dot products. This rewriting consequently does not by itself establish
backward stability. The pivot dependence is made explicit instead in \S\ref{sec:conditional}.
\end{remark}

The permutation $P_j$ must be carried explicitly, since omitting it from the accumulated product of
\S\ref{sec:accumulation} produces a quantity that is not the exactly applied tableau at all,
and the resulting spurious defect, verified directly, is larger by up to fourteen orders
of magnitude on well conditioned data.

\subsection{Accumulation and amplification}
\label{sec:accumulation}

Write $\widehat{\mathcal{B}}_{i:k} = \prod_{j=i}^{k}P_j\widehat B_j$, with
$\widehat{\mathcal{B}}_{k+1:k}=I$, and let $\widehat V_1$ be the starting tableau
($\widehat V_1=I$ in~\cite{CastilloCJPC2000}).

\begin{theorem}[Accumulation]\label{thm:accum}
After $N=n$ steps,
%
\begin{equation}\label{eq:accum}
\widehat V_{N+1}
=
\widetilde V_{N+1} + E_{N+1},
\qquad
\widetilde V_{N+1} = \widehat V_1\widehat{\mathcal{B}}_{1:N},
\qquad
E_{N+1} = \sum_{j=1}^{N} F_j\,\widehat{\mathcal{B}}_{j+1:N},
\end{equation}
%
and consequently
%
\begin{equation}\label{eq:accum-bound}
\bigl\|E_{N+1}\bigr\|
\;\le\;
\gamma_2 \sum_{j=1}^{N}
\bigl\|\,\bigl|\widehat V_jP_j\bigr|\,\bigl|\widehat B_j\bigr|\,\bigr\|
\;\bigl\|\widehat{\mathcal{B}}_{j+1:N}\bigr\| .
\end{equation}
\end{theorem}

\begin{proof}
Induction on $j$ using~\eqref{eq:onestep}, followed by the triangle inequality and
submultiplicativity.
\end{proof}

Theorem~\ref{thm:accum} is the analogue, for the normalized inverse-forming variant, of the
analysis Fletcher~\cite[Thm.~1]{Fletcher1997} gives for the LIU factorization, and we claim no
priority for it; his estimate is stated for the triangular factors and involves the growth of
the $\mathbb{L}^{(p)}$, whereas~\eqref{eq:accum-bound} concerns the explicitly formed tableau.
What we add is the observation of \S\ref{sec:defect} that~\eqref{eq:accum} does not by itself
bound the backward error of the inverse, and the finding of \S\ref{sec:law} that it is far from
sharp.

\subsection{Backward error of the computed inverse}
\label{sec:defect}

Define the residual and the \emph{transformation defect}
%
\begin{equation}\label{eq:defect}
R_N = I - A\widehat V_{N+1},
\qquad
D_N = I - A\widetilde V_{N+1} .
\end{equation}
%
The defect measures the extent to which the computed elementary transformations, applied
\emph{exactly}, fail to invert $A$. It need not be zero: the coefficients $\hat\alpha_k$ and
$\hat\beta$ are formed from rounded inner products and rounded divisions, so in general
$\widetilde V_{N+1}$ is the exact product of factors that differ from their exact-arithmetic
counterparts.

\begin{theorem}[Backward error of the inverse]\label{thm:backward}
Let $\widehat V_{N+1}$ be nonsingular. Then
%
\begin{equation}\label{eq:decomposition}
R_N = D_N - A\,E_{N+1},
\end{equation}
%
and the unique matrix $\Delta A$ with $(A+\Delta A)\widehat V_{N+1}=I$ is
%
\begin{equation}\label{eq:deltaA}
\Delta A = R_N\,\widehat V_{N+1}^{-1},
\end{equation}
%
so that the normwise inverse backward error satisfies
%
\begin{equation}\label{eq:etainv-bound}
\eta_{\mathrm{inv}} := \frac{\|\Delta A\|}{\|A\|}
\;\le\;
\left(\frac{\|D_N\|}{\|A\|} + \|E_{N+1}\|\right)\bigl\|\widehat V_{N+1}^{-1}\bigr\| .
\end{equation}
\end{theorem}

\begin{proof}
Substituting $\widehat V_{N+1}=\widetilde V_{N+1}+E_{N+1}$ into the definition of $R_N$
gives $R_N = I-A\widetilde V_{N+1}-AE_{N+1} = D_N - AE_{N+1}$, which
is~\eqref{eq:decomposition}. Any $\Delta A$ with $(A+\Delta A)\widehat V_{N+1}=I$ satisfies
$\Delta A\widehat V_{N+1}=R_N$, and nonsingularity of $\widehat V_{N+1}$ gives uniqueness
and~\eqref{eq:deltaA}. The bound follows by the triangle inequality applied
to~\eqref{eq:decomposition} together with~\eqref{eq:accum-bound}.
\end{proof}

Two features deserve comment. First, \eqref{eq:decomposition} is an exact identity, not an
estimate, and it is the structural content of the theorem: the residual contains two distinct
algebraic contributions. They are not statistically or dynamically independent---indeed
Theorem~\ref{thm:defect-structure} contains an explicit coupling through the previously accumulated
application error---but a bound on $E_{N+1}$ alone is nevertheless \emph{not} a bound on
$\eta_{\mathrm{inv}}$. Setting $D_N=0$ is legitimate only under the additional hypothesis that
the computed transformations invert $A$ exactly when applied exactly, which~\eqref{eq:dots}
shows to be false in general. Second, the two contributions are not comparably informative in the
campaign: the defect contribution tracks $\eta_{\mathrm{inv}}$ to within half an order of magnitude
in the median, whereas the full right-hand side of~\eqref{eq:etainv-bound} exceeds it by a median
factor of $2.3\times10^{5}$ (\S\ref{sec:law}).

\subsection{Structure of the transformation defect}
\label{sec:defect-structure}

The defect was defined in \S\ref{sec:defect} by what it measures; we now show how it is
generated. An exact telescoped representation assembles $D_N$, row by row, from a local
\emph{formation defect} created at the step in which each row is biorthogonalized, whose
componentwise bound exhibits the scaled pivot and the normalized absolute sensitivity of the dot
products.

\begin{proposition}[Local formation defect]\label{prop:local-defect}
Assume $(n+1)u<1$ and the standard floating-point model without underflow or overflow. At a
successful step $j$ with nonzero computed pivot, let
$\overline V_j=\widehat V_jP_j$, $\overline v_k^{\,j}=\overline V_je_k$, and
$s^T=a_j^T\overline V_j$ be the exact dot products with the computed tableau in permuted column
order. Define the formation defect $\varphi_j^T=s^T\widehat B_j-e_j^T$. Then
%
\begin{equation}\label{eq:phi-components}
\varphi_{j,j}=s_j\hat\beta-1,
\qquad
\varphi_{j,k}=s_k-s_j\hat\alpha_k \quad (k\ne j),
\end{equation}
%
and
%
\begin{equation}\label{eq:phi-pair}
|\varphi_{j,j}|
\le u+\gamma_n|a_j|^T|\overline v_j^{\,j}|\,|\hat\beta|,
\qquad
|\varphi_{j,k}|
\le \gamma_{n+1}|a_j|^T|\overline v_k^{\,j}|
+\gamma_n|a_j|^T|\overline v_j^{\,j}|\,|\hat\alpha_k|.
\end{equation}
%
These inequalities imply the compact componentwise bound
%
\begin{equation}\label{eq:phi-compact}
|\varphi_j^T|
\le
\gamma_{n+1}|a_j|^T|\widehat V_jP_j|\,|\widehat B_j|.
\end{equation}
\end{proposition}

\begin{proof}
Write $\overline t_k^{\,j}=s_k+\delta_k$, where
$|\delta_k|\le\gamma_n|a_j|^T|\overline v_k^{\,j}|$, and write the rounded quotient and reciprocal
as
$\hat\alpha_k=(\overline t_k^{\,j}/\overline t_j^{\,j})(1+\theta_k)$ and
$\hat\beta=(1+\theta_0)/\overline t_j^{\,j}$, with
$|\theta_k|,|\theta_0|\le u$.  Substitution gives the component identities and the two estimates
in~\eqref{eq:phi-pair}.  For the pivot component put
$q_j=|a_j|^T|\overline v_j^{\,j}|$.  Since
$|\overline t_j^{\,j}|\le(1+\gamma_n)q_j$ and
$|\hat\beta|\ge(1-u)/|\overline t_j^{\,j}|$,
$q_j|\hat\beta|\ge(1-u)(1-nu)$.  Hence
\[
\gamma_{n+1}-\gamma_n-
\frac{u}{(1-u)(1-nu)}
=
\frac{nu^2}{(1-u)(1-nu)(1-(n+1)u)}\ge0,
\]
which absorbs the isolated reciprocal-rounding term into $\gamma_{n+1}$.  The off-pivot
components then follow directly from the columns of $|\widehat B_j|$.  SM2 gives the componentwise
calculation in full.
\end{proof}

\begin{theorem}[Telescoped defect]\label{thm:defect-structure}
With $G_j = e_j\varphi_j^T$ and
$E_j = \widehat V_j - \widehat V_1\widehat{\mathcal B}_{1:j-1}$ the application error
accumulated before step $j$ (so $E_1 = 0$),
%
\begin{equation}\label{eq:defect-telescope}
D_N
\;=\;
-\sum_{j=1}^{N} G_j\,\widehat{\mathcal{B}}_{j+1:N}
\;+\;
\sum_{j=1}^{N} e_j\,a_j^T E_j\,\widehat{\mathcal{B}}_{j:N} .
\end{equation}
\end{theorem}

\begin{proof}
Write $W_j = \widehat V_1\widehat{\mathcal B}_{1:j-1}$ and, for $i\le j$, let
$g_i^{(j)T} = a_i^TW_{j+1} - e_i^T$. For $l>i$ the permutation $P_l$ fixes coordinate $i$, since
$p_l\ge l>i$, and $e_i^T\widehat B_l = e_i^T$, so $g_i^{(l)T} = g_i^{(l-1)T}P_l\widehat B_l$: the
defect of a processed row propagates multiplicatively and no new term is created. Hence
$g_i^{(N)T} = g_i^{(i)T}\widehat{\mathcal B}_{i+1:N}$, while at step $i$ itself
$g_i^{(i)T} = \varphi_i^T - a_i^TE_iP_i\widehat B_i$ by $W_i=\widehat V_i-E_i$. Row $i$ of
$D_N=I-\Pi A\widetilde V_{N+1}$ is $-g_i^{(N)T}$, and summing over $i$
gives~\eqref{eq:defect-telescope}.
\end{proof}

The representation~\eqref{eq:defect-telescope} converts the defect from a defined quantity
into an analyzed one. The first sum is created by \emph{forming} the transformations from rounded
dot products; the second couples those formation rows to the application errors already
accumulated, vanishes when the transformations are applied exactly, and is controlled by the same
growth quantities as~\eqref{eq:accum-bound} (SM2).

To interpret the local bound, for a nonzero column vector $v$ define the normalized absolute
dot-product sensitivity
%
\begin{equation}\label{eq:chi}
\chi(a_j,v)
=
\frac{|a_j|^T|v|}{\|a_j\|_2\,\|v\|_2}
\in[0,1].
\end{equation}
%
This is not the relative condition number of the dot product. For a nonzero exact product
$a_j^Tv$ the latter is $|a_j|^T|v|/|a_j^Tv|$ and can be arbitrarily large. After normalization,
the pivot component of~\eqref{eq:phi-pair} is governed by a factor of the form
$\gamma_n\chi(a_j,\overline v_j^{\,j})/\hat p_j^{\rm sc}$, while the off-pivot components carry
analogous factors weighted by $|\hat\alpha_k|$. Thus a small scaled pivot magnifies the formation
defect through $|\hat\beta|$ and the multipliers. Dot-product cancellation supplies an additional mechanism not
captured by the pivot alone; this is precisely where the empirical pivot diagnostic is weakest on
the Vandermonde family in \S\ref{sec:predictors}. The propagated trailing products remain the
non-sharp part of the resulting global estimate, just as in \S\ref{sec:accumulation}.

\subsection{A growth factor for the transformation}
\label{sec:growth}

The bound~\eqref{eq:accum-bound} and the componentwise estimate of
Proposition~\ref{prop:local-defect} control the residual through the same combination of computed
quantities, which it is convenient to isolate. Throughout this subsection norms are the
$\infty$-norm, for which $\||X|\|_\infty=\|X\|_\infty$ so that the componentwise absolute values
in~\eqref{eq:accum-bound} cost nothing; we write
%
\begin{equation}\label{eq:rinv-inf}
r_R^{\infty} = \frac{\bigl\|I-\Pi A\widehat V_{N+1}\bigr\|_\infty}
                    {\|A\|_\infty\,\|\widehat V_{N+1}\|_\infty}
\end{equation}
%
for the $\infty$-norm counterpart of the scaled right residual~\eqref{eq:rinv}. Standard
norm equivalence gives
$ n^{-3/2}r_R\le r_R^{\infty}\le n^{3/2}r_R$; no sharper equivalence is assumed below.

\begin{definition}\label{def:growth}
The \emph{growth factor} of the pivoting transformation is
%
\begin{equation}\label{eq:growth}
\Gamma_N
=
\frac{1}{\|\widehat V_{N+1}\|_\infty}
\sum_{l=1}^{N}
\bigl\|\,\bigl|\widehat V_lP_l\bigr|\,\bigl|\widehat B_l\bigr|\,\bigr\|_\infty
\;\bigl\|\widehat{\mathcal{B}}_{l+1:N}\bigr\|_\infty .
\end{equation}
\end{definition}

When an exact-arithmetic trajectory is discussed below, the same symbol $\Gamma_N$ denotes the
corresponding expression with the hats omitted. This convention keeps the algebraic structure
visible while distinguishing exact-arithmetic bounds from the computed residual estimate.

\begin{proposition}[Residual bound]\label{prop:growth-bound}
$r_R^{\infty} \le \gamma_{n+3}\,\Gamma_N$.
\end{proposition}

\begin{proof}
Theorem~\ref{thm:defect-structure} and the row-wise argument of its proof give
$R_N = -\sum_{i} e_i\varphi_i^T\widehat{\mathcal{B}}_{i+1:N}
- \sum_{l} \Pi_l A F_l\widehat{\mathcal{B}}_{l+1:N}$ with $\Pi_l=\sum_{i\le l}e_ie_{\pi_i}^T$.
For a row vector, $\|v^TM\|_1\le\|v\|_1\|M\|_\infty$; hence by
Proposition~\ref{prop:local-defect} and $\|a_i\|_1\le\|A\|_\infty$ the $i$-th row of the first
sum has $1$-norm at most $\gamma_{n+1}\|A\|_\infty\||\widehat V_iP_i||\widehat B_i|\|_\infty
\|\widehat{\mathcal{B}}_{i+1:N}\|_\infty$, and the $\infty$-norm of the sum is bounded by the
sum of these. The second sum is bounded termwise using $\|\Pi_l\|_\infty\le1$ and
$|F_l|\le\gamma_2|\widehat V_lP_l||\widehat B_l|$ from Lemma~\ref{lem:onestep}. Adding, with
$\gamma_{n+1}+\gamma_2\le\gamma_{n+3}$, gives the claim.
\end{proof}

\begin{proposition}[A linear floor for the growth factor]\label{prop:growth-floor}
The computed growth factor satisfies
%
\begin{equation}\label{eq:growth-floor}
\Gamma_N \ge \frac{N}{1+N\gamma_2}.
\end{equation}
%
For fixed $N$ this is $\Gamma_N\ge N-2N^2u+O(u^2)$ as $u\to0$.
\end{proposition}

\begin{proof}
Fix $l$. Submultiplicativity and the entrywise inequality $|X||Y|\ge|XY|$ give
$\||\widehat V_lP_l||\widehat B_l|\|_\infty\|\widehat{\mathcal{B}}_{l+1:N}\|_\infty
\ge\|\widehat V_lP_l\widehat B_l\widehat{\mathcal{B}}_{l+1:N}\|_\infty$. Iterating
\eqref{eq:onestep} from step $l$ onwards gives the exact identity
%
\[
\widehat V_lP_l\widehat B_l\widehat{\mathcal B}_{l+1:N}
=
\widehat V_{N+1}-\sum_{m=l}^{N}F_m\widehat{\mathcal B}_{m+1:N}.
\]
%
The sum on the right has infinity norm at most
$\gamma_2\|\widehat V_{N+1}\|_\infty\Gamma_N$ by~\eqref{eq:accum-bound} and
\eqref{eq:growth}. After division by $\|\widehat V_{N+1}\|_\infty$, each of the $N$ summands
of~\eqref{eq:growth} is therefore at least $1-\gamma_2\Gamma_N$. Hence
$\Gamma_N\ge N(1-\gamma_2\Gamma_N)$, which is equivalent to~\eqref{eq:growth-floor}.
\end{proof}

\begin{corollary}\label{cor:floor}
The right-hand side of Proposition~\textup{\ref{prop:growth-bound}} itself has the finite-$u$
floor
%
\[
\gamma_{n+3}\Gamma_N
\ge
\gamma_{n+3}\frac{n}{1+n\gamma_2}.
\]
%
Consequently an estimate of this form cannot certify a scaled residual below
$n(n+3)u+O(u^2)$ at first order, for fixed $n$, whatever the matrix or pivot sequence.
\end{corollary}

Corollary~\ref{cor:floor} isolates a part of the pessimism of the estimate that is structural
rather than attributable to additional tableau growth: bounding separately two factors whose
product telescopes exactly incurs a contribution of order $N$ once per step even when the
transformation is otherwise well behaved, and the measured residuals of \S\ref{sec:experiments}
lie far below this first-order floor.

\subsection{The intermediate tableaux}
\label{sec:tableaux}

The growth factor is governed by the tableaux, which for the unpivoted process have a closed
form. For $1\le j\le n$ write $A_j=A(1{:}j,1{:}j)$ for the leading principal submatrix and
$A_j'=A(1{:}j,j{+}1{:}n)$ for the remaining block of the same rows.

\begin{lemma}[Explicit tableaux]\label{lem:tableaux}
Apply~\eqref{eq:castillo-pivot} in exact arithmetic with $V^1=I$, the rows of $A$ in natural
order, and no column interchanges, and suppose $A_1,\dots,A_j$ are nonsingular. Then
%
\begin{equation}\label{eq:tableau-explicit}
V^{j+1}
=
\begin{bmatrix} A_j^{-1} & -A_j^{-1}A_j' \\[2pt] 0 & I_{n-j}\end{bmatrix},
\qquad
\bigl(V^{j+1}\bigr)^{-1}
=
\begin{bmatrix} A_j & A_j' \\[2pt] 0 & I_{n-j}\end{bmatrix}.
\end{equation}
\end{lemma}

\begin{proof}
Rows $m>j$ remain the corresponding rows of the identity.  The invariant
$[A_j\ A_j']V^{j+1}=[I_j\ 0]$ then gives $A_jX=I_j$ and $A_jY+A_j'=0$, yielding
\eqref{eq:tableau-explicit}; direct multiplication gives the inverse.
\end{proof}

Lemma~\ref{lem:tableaux} makes precise the warning of Hestenes recalled in
\S\ref{sec:genealogy}: started from $V^{(0)}=I$ the tableaux are the inverses of the leading
principal submatrices of $A$ itself, bordered by the corresponding off-diagonal block, so the
process inherits whatever ill-conditioning those submatrices possess. Since
$\|(V^{j+1})^{-1}\|_\infty\le\max\{1,\|A\|_\infty\}$ uniformly in $j$, the trailing products
in~\eqref{eq:growth} are controlled by the data and all the growth resides in the tableaux.

\subsection{Matrices with diagonal dominance}
\label{sec:dd}

For the classes in which the growth factor of Gaussian elimination is bounded by classical
theorems, Lemma~\ref{lem:tableaux} transfers those bounds to $\Gamma_N$. We carry out the
argument for strict diagonal dominance, where the constants can be made explicit.

Let $A$ be strictly diagonally dominant by rows,
%
\begin{equation}\label{eq:dd-gap}
\alpha := \min_{1\le i\le n}\Bigl(|a_{ii}|-\sigma_i\Bigr) > 0,
\qquad
\sigma_i := \sum_{k\ne i}|a_{ik}| ,
\end{equation}
%
and normalize $\|A\|_\infty=1$, so that $0<\alpha\le1$ and $\sigma_i\le1-\alpha$.

\begin{theorem}[Growth for diagonally dominant matrices]\label{thm:dd}
Let $A$ satisfy~\eqref{eq:dd-gap} with $\|A\|_\infty=1$, and apply
\eqref{eq:castillo-pivot} in exact arithmetic with $V^1=I$ and the rows of $A$ in natural order.
Then:
\begin{enumerate}
\item every leading principal submatrix $A_j$ satisfies~\eqref{eq:dd-gap} with the same $\alpha$,
and $\|A_j^{-1}\|_\infty\le1/\alpha$;
\item every Schur complement is strictly row-diagonally dominant with gap at least $\alpha$, and
every pivot satisfies $|t_l^{\,l}|\ge\alpha$;
\item for every active row and every $k>l$,
$|t_l^{\,l}|-|t_k^{\,l}|\ge\alpha$; hence the first-nonzero and maximum-modulus column rules both
select the diagonal entry at every step and follow the same path without column interchanges;
\item $\|V^j\|_\infty\le2/\alpha$ and $\|(V^j)^{-1}\|_\infty\le1$ for every $j$, while
$\|B_l\|_\infty\le3/\alpha^2$ for every $l$;
\item the exact-arithmetic growth factor satisfies
%
\begin{equation}\label{eq:dd-growth}
\Gamma_N\le\frac{6n}{\alpha^3}.
\end{equation}
\end{enumerate}
\end{theorem}

\begin{proof}[Proof sketch; details in SM2]
Every $A_j$ inherits the gap $\alpha$, and the standard maximum-component argument gives
$\|A_j^{-1}\|_\infty\le1/\alpha$.  One elimination step preserves that gap because
\[
|a'_{ii}|-\sum_{k\ge2,\,k\ne i}|a'_{ik}|
\ge (|a_{ii}|-\sigma_i)
+|a_{i1}|\left(1-\frac{\sigma_1}{|a_{11}|}\right)\ge\alpha.
\]
By Lemma~\ref{lem:tableaux}, the tail $(t_l^{\,l},\ldots,t_n^{\,l})$ is the active Schur row,
so the diagonal pivot exceeds every competitor in modulus by at least $\alpha$.  The same lemma
gives $\|V^j\|_\infty\le2/\alpha$ and $\|(V^j)^{-1}\|_\infty\le1$; consequently
$\|B_l\|_\infty\le3/\alpha^2$.  Finally
$\mathcal B_{l+1:N}=(V^{l+1})^{-1}V^{N+1}$, so the terminal tableau norm cancels in
\eqref{eq:growth}, giving $\Gamma_N\le6n/\alpha^3$.
\end{proof}

\begin{corollary}[First-order residual bound on the diagonal-dominance path]\label{cor:dd-fp}
Fix $A$ and $n$ as in Theorem~\ref{thm:dd}. In the standard floating-point model, excluding
underflow and overflow, for all sufficiently small $u$ the computed first-nonzero and
maximum-modulus variants follow the exact pivot path. Along that path,
%
\begin{equation}\label{eq:dd-residual}
r_R^{\infty}
\le
\frac{6n(n+3)}{\alpha^3}\,u+O(u^2),
\qquad u\to0,
\end{equation}
%
where the constant hidden in $O(u^2)$ may depend on the fixed matrix and on $n$.
\end{corollary}

\begin{proof}
The proof of Theorem~\ref{thm:dd} gives more than strict comparison: in every exact active Schur
complement,
$|t_l^{\,l}|\ge\alpha$ and
$|t_l^{\,l}|-|t_k^{\,l}|\ge\alpha$ for each $k>l$, because the full row-dominance gap is at least
$\alpha$. For fixed $A$ and $n$, a finite induction along the exact path gives
$\widehat V_j=V^j+O(u)$ and computed dot products equal to their exact-path counterparts plus
$O(u)$. Hence, once $u$ is sufficiently small, the active computed pivot remains nonzero and the
ordering of the computed moduli cannot change; the first-nonzero rule also accepts the active
entry. The same induction then gives $\widehat B_j=B_j+O(u)$ at every step. Consequently the
computed growth factor is $\Gamma_N+O(u)$. Proposition~\ref{prop:growth-bound},
$\gamma_{n+3}=(n+3)u+O(u^2)$, and~\eqref{eq:dd-growth} give~\eqref{eq:dd-residual}.
\end{proof}

\begin{remark}\label{rem:dd-caveats}
The distinction between Theorem~\ref{thm:dd} and Corollary~\ref{cor:dd-fp} is deliberate: the
former is an exact-arithmetic, matrix-class bound, whereas the latter is a first-order statement
about the computed trajectory. The maximum-scaled column rule is not covered because its
normalizing denominators $\|v_k\|_2$ are not entries of the Schur complement; indeed
\S\ref{sec:dd-numerics} exhibits interchanges under that rule. Row pivoting is also outside the
theorem because it replaces the leading block of Lemma~\ref{lem:tableaux} by
$A(\pi(1{:}j),1{:}j)$. Finally, the exponent $3$ in~\eqref{eq:dd-growth} is not claimed sharp. It
arises once from $\|V^j\|_\infty$, once from the reciprocal pivot, and once from bounding the full
row $t$ by $\|V^j\|_\infty\|a_j\|_1$; the numerical evidence in \S\ref{sec:dd-numerics} shows that
this last chain of separate estimates is highly pessimistic.
\end{remark}

\subsection{Backward error of an individual linear system}
\label{sec:solution-theory}

If the computed inverse is used to solve $Ax=b$ by forming
$\widehat x=\fl(\widehat V_{N+1}b)$, the resulting residual depends on the action of $R_N$ on the
particular right-hand side, not only on its norm.

\begin{proposition}\label{prop:sol-vs-inv}
Let $b\ne0$ and $\widehat x=\fl(\widehat V_{N+1}b)$. There exists $\Delta V$ with
$|\Delta V|\le\gamma_n|\widehat V_{N+1}|$ such that
$\widehat x=(\widehat V_{N+1}+\Delta V)b$, and
%
\begin{equation}\label{eq:sol-residual}
b-A\widehat x=R_Nb-A\Delta Vb.
\end{equation}
%
Consequently,
%
\begin{equation}\label{eq:sol-bound}
\eta_{\mathrm{sol}}
\le
\frac{\|R_Nb\|_2}{\|b\|_2}
+
\gamma_n\|A\|_2\,\bigl\|\,|\widehat V_{N+1}|\,\bigr\|_2,
\end{equation}
%
and, using $R_N=\Delta A\widehat V_{N+1}$ from Theorem~\ref{thm:backward},
%
\begin{equation}\label{eq:sol-inv-bound}
\eta_{\mathrm{sol}}
\le
\eta_{\mathrm{inv}}\,\|A\|_2\,\|\widehat V_{N+1}\|_2
+
\gamma_n\|A\|_2\,\bigl\|\,|\widehat V_{N+1}|\,\bigr\|_2.
\end{equation}
\end{proposition}

The proof is given in SM2. Equation~\eqref{eq:sol-residual} shows why a particular right-hand side
can have small solution backward error even when the matrix inverse has a large backward error.
Conversely,~\eqref{eq:sol-inv-bound} makes explicit the factor multiplying
$\eta_{\mathrm{inv}}$ in a uniform-in-$b$ statement; when
$\widehat V_{N+1}\approx A^{-1}$ this factor is of order $\kappa_2(A)$. This is also why
comparisons based on $\eta_{\mathrm{sol}}$ require the underlying problem to be resolved in the
working precision (\S\ref{sec:scaling}). Iterative refinement changes $\widehat x$ and therefore
$\eta_{\mathrm{sol}}$, but it does not alter the stored inverse or $R_N$.

\subsection{A conditional statement, and the role of the pivot}
\label{sec:conditional}

The elementary transformation~\eqref{eq:Bhat} has
$\|\widehat B_j\|\ge|\hat\beta|$ and, because
$\hat\beta=\fl(1/\overline t_j^{\,j})$,
%
\[
|\hat\beta|\ge\frac{1-u}{|\overline t_j^{\,j}|}.
\]
%
Thus the pivot enters the local amplification directly. We measure it by the scale-invariant
quantity
%
\begin{equation}\label{eq:pmin}
p^{\rm sc}_{\min}=\min_{1\le j\le N}\hat p_j^{\rm sc},
\qquad
\hat p_j^{\rm sc}
=
\frac{|\overline t_j^{\,j}|}
     {\|a_j\|_2\,\|\overline v_j^{\,j}\|_2}.
\end{equation}
%
Equivalently, $\overline v_j^{\,j}=\widehat v_{p_j}^{\,j}$ and
$\overline t_j^{\,j}=\hat t_{p_j}^{\,j}$ in the pre-exchange indexing. Equation~\eqref{eq:dots}
gives the normalized absolute perturbation bound
%
\[
\left|
\frac{\overline t_j^{\,j}-a_j^T\overline v_j^{\,j}}
     {\|a_j\|_2\,\|\overline v_j^{\,j}\|_2}
\right|
\le\gamma_n\chi(a_j,\overline v_j^{\,j});
\]
%
this is not a relative-error bound for the dot product when cancellation is severe.

Theorem~\ref{thm:backward} immediately gives the following conditional statement. If along the
computed sequence
%
\begin{equation}\label{eq:hypotheses}
\frac{\|D_N\|}{\|A\|}\le c_Du,
\qquad
\|E_{N+1}\|\le c_Eu,
\qquad
\|\widehat V_{N+1}^{-1}\|\le c_V,
\end{equation}
%
then
$\eta_{\mathrm{inv}}\le(c_D+c_E)c_Vu$. The mathematical issue is therefore not this immediate
implication but which hypotheses can be controlled for the algorithm.

\begin{observation}[Local pivot diagnostic]\label{obs:pivot}
Over the reliability-screened records of the mechanistic campaign, the empirical inequality
%
\begin{equation}\label{eq:defect-pivot}
\frac{\|D_N\|_2}{\|A\|_2}
\le
\frac{un}{p^{\rm sc}_{\min}}
\end{equation}
%
holds in $88.7\%$ of cases overall, in $99.7\%$ for \texttt{R1\_C1}, and in $99.96\%$ for
\texttt{R2\_C2}; the analogous inequality for $\eta_{\mathrm{inv}}$ holds in $99.0\%$ of the
screened cases. The defect itself is a tighter ex-post diagnostic of $\eta_{\mathrm{inv}}$ under
the ratio-spread criterion of \S\ref{sec:predictors}; the significance of
$p^{\rm sc}_{\min}$ is instead that it is a local trajectory quantity and is the strongest such
indicator among those tested. The relation is incomplete, with Vandermonde matrices providing the
principal failures. Independently, \S\ref{sec:adversarial} proves that $\kappa_2(A)$ alone cannot
control the intermediate tableau growth: an orthogonal family with $\kappa_2(A)=1$ has
$p^{\rm sc}_{\min}=\varepsilon$ and growth asymptotic to $\sqrt2/\varepsilon$ under the
first-nonzero rule.
\end{observation}

Observation~\ref{obs:pivot} is empirical, not a theorem. It identifies the local mechanism that a
sharper conditional analysis would have to propagate without the severe loss introduced by the
growth products. The regression evidence in SM6 also shows that growth variables retain
class-dependent information out of sample; the conclusion is therefore not that growth is
irrelevant, but that the classical growth estimate is too pessimistic and that the scaled pivot is
the dominant local diagnostic of defect formation.

\section{Numerical campaign}
\label{sec:experiments}

Two campaigns are reported. The \emph{mechanistic} campaign tests the quantities appearing in
\eqref{eq:hypotheses} and identifies which ones are informative for the observed backward error; it
comprises $18\,744$ matrix/right-hand-side configurations over seven families and $552\,747$
method/refinement records. The \emph{large-scale} campaign tests whether those conclusions survive
dimension, structure and massive sampling, and measures quantities that the theory defines but
does not sharply bound. It contains $14.7$ million complete matrices over fourteen families,
thirteen dimensions $8\le n\le512$, and $1.47\times10^8$ inverse records. The $20\,000$-sample
core is complete through $n=181$; dimensions $256$, $362$, and $512$ are analyzed through a
balanced extension of $2000$ realizations per stochastic cell recovered after the wall-clock limit.
Because the recovery follows the deterministic task/manifest ordering, no missing-at-random claim
is made for that extension. Design, generation parameters, completeness accounting and provenance
are in SM3 and SM7.

\subsection{Design and instrumentation}
\label{sec:design}

Both campaigns are evaluated in binary32 and binary64, with $u_{32}=2^{-24}$ and
$u_{64}=2^{-53}$, and with zero, one and two steps of iterative refinement. A single instrumented
implementation records the elementary transformations, the permutations, the tableaux and the
quantities of \S\ref{sec:growth}, so every diagnostic below derives from the same computed data
as the inverse itself.

Two pivoting decisions are available in~\eqref{eq:castillo-pivot} and they are logically
independent: the \emph{column rule} selects the pivot column at step $j$, and the \emph{row
rule} selects which of the remaining rows of $A$ is processed. The column search is necessarily
restricted to columns not previously used as pivots, since a column $v_i$ with $a_i^Tv_i=1$
cannot be reused without destroying the biorthogonality already
established~\cite[Algorithm~1, Step~3]{CastilloPSM2006}. The row rule carries no such
restriction and is the axis along which the published regimes differ: G\'ati~\cite{Gati2013}
uses no pivoting of either kind, whereas Fletcher~\cite[\S2]{Fletcher1997} permutes rows to
bring the largest $|\mathbf{l}_j^Ta_k|$ to the diagonal, exactly equivalent to partial pivoting
in Gaussian elimination. We therefore consider the five combinations of
Table~\ref{tab:variants}, the scaled score being
$s_{jk} = |\hat t^{\,j}_k|/(\|a_j\|_2\|\widehat v_k^{\,j}\|_2)$, evaluated from the computed dot
products consistently with~\eqref{eq:pmin}. The references are LU with partial pivoting,
Householder QR, the SVD, and Gauss--Jordan inversion with partial and with complete pivoting.

\begin{table}
\centering
\caption{Pivoting variants. \texttt{R0\_C0} is the rule of~\cite{CastilloCJPC2000};
\texttt{R1\_C1} is the closest available analogue of the LIU factorization
of~\cite{Fletcher1997} with partial pivoting.}
\label{tab:variants}
\begin{tabular}{lll}
\hline
Variant & Row rule & Column rule \\
\hline
\texttt{R0\_C0} & rows in order & first numerically nonzero pivot \\
\texttt{R0\_C1} & rows in order & maximum absolute pivot \\
\texttt{R0\_C2} & rows in order & maximum scaled pivot \\
\texttt{R1\_C1} & maximum absolute & maximum absolute pivot \\
\texttt{R2\_C2} & maximum scaled  & maximum scaled pivot \\
\hline
\end{tabular}
\end{table}

In Table~\ref{tab:variants}, ``first numerically nonzero'' means the first admissible computed
pivot that is unequal to zero in the working format; no separate tolerance-based acceptance rule is
part of this variant. The first-order path statement of Corollary~\ref{cor:dd-fp} refers to this
rule; a tolerance-based modification would require its threshold to be specified as $u\to0$.

\subsection{Measured quantities and their reliability}
\label{sec:reliability}

The exact normwise inverse backward error of Theorem~\ref{thm:backward} is
%
\begin{equation}\label{eq:etainv}
\eta_{\mathrm{inv}}
=
\frac{\|(I-A\widehat V_{N+1})\widehat V_{N+1}^{-1}\|_2}{\|A\|_2}.
\end{equation}
%
The campaign evaluates a floating-point estimator of this exact metric by solving with the
computed tableau. When $\widehat V_{N+1}$ is itself unresolved in the working precision, that
solve is the operation whose reliability is in question. We therefore use the same conservative conditioning screen for every
variant. Let $W$ denote the separately maintained inverse of the accumulated product of the
computed transformations, cast with $\widehat V_{N+1}$ to binary64 for the screen, and define
%
\begin{equation}\label{eq:reliability-delta}
\delta_F
=
\bigl\|I-\fl_{64}(W\widehat V_{N+1})\bigr\|_F
+
\gamma_n^{(64)}\,\|W\|_F\,\|\widehat V_{N+1}\|_F,
\end{equation}
%
where $\gamma_n^{(64)}$ uses binary64 unit roundoff. The motivation is the standard product
error bound: if the two displayed terms were evaluated exactly, then
%
\[
\|I-W\widehat V_{N+1}\|_2
\le
\|I-W\widehat V_{N+1}\|_F
\le\delta_F.
\]
%
Thus $\delta_F<1$ implies
$\widehat V_{N+1}^{-1}=(I-H)^{-1}W$ with
$H=I-W\widehat V_{N+1}$ and hence
%
\[
\|\widehat V_{N+1}^{-1}\|_2
\le\frac{\|W\|_2}{1-\delta_F},
\qquad
\kappa_2(\widehat V_{N+1})
\le
\frac{\|\widehat V_{N+1}\|_F\,\|W\|_F}{1-\delta_F}.
\]
%
This motivates the implemented screen
%
\begin{equation}\label{eq:reliability}
\kappa^{\rm scr}
=
\frac{\|\widehat V_{N+1}\|_F\,\|W\|_F}{1-\delta_F}
\end{equation}
%
and the requirement $u\kappa^{\rm scr}<10^{-2}$ before the estimator of
$\eta_{\mathrm{inv}}$ is reported. We call these records \emph{reliability-screened}. The
terminology is intentionally weaker than ``exactly certified'': the norms and scalar combinations
in the screen are themselves evaluated in binary64 without directed rounding, so
\eqref{eq:reliability-delta} is an operational conservative screen rather than a formal interval
certificate or a direct high-precision evaluation of $\kappa_2(\widehat V_{N+1})$. The threshold
was fixed before the campaign and applied identically to all variants.

In the mechanistic campaign $52\,184$ of the $91\,466$ successful unrefined Castillo records
($57.1\%$) pass the screen. In the large-scale campaign the corresponding rate is $48\%$ in
binary32 and $81\%$ in binary64, with less than one percentage point of variation among pivot
variants. Comparisons based on $\eta_{\mathrm{inv}}$ are therefore reported only on the screened
subset with its size stated. The scaled right inverse residual
%
\begin{equation}\label{eq:rinv}
r_R
=
\frac{\|I-A\widehat V_{N+1}\|_2}{\|A\|_2\,\|\widehat V_{N+1}\|_2}
\end{equation}
%
requires no solve and is available for every successful record. SM4 gives the full reliability
accounting and the exact implementation of the screen.

The instrumentation was checked independently. A verifier reproduced the working-precision pivot
path of the $91\,466$ mechanistic inverse records and evaluated both~\eqref{eq:decomposition} and
the row-wise reconstruction of $D_N$ from Theorem~\ref{thm:defect-structure} in variable precision,
with precision and tolerance scaled to the observed trailing-product growth. Both identities passed
on all $91\,219$ records for which the pivot path was reproduced. Since the identities themselves
are algebraic, this validates the instrumentation and the treatment of row and column permutations,
not the theorems by experiment (SM11).

\subsection{Comparison of pivoting rules}
\label{sec:comparison}

\begin{table}
\centering
\caption{Scaled right inverse residual~\eqref{eq:rinv} over all successful cases of the
mechanistic campaign, and reliability-screened inverse backward error~\eqref{eq:etainv}. LU, QR and the SVD
do not form an explicit inverse and are not comparable here.}
\label{tab:comparison}
\footnotesize
\setlength{\tabcolsep}{2.5pt}
\begin{tabular}{lrrrrrr}
\hline
& \multicolumn{3}{c}{$r_R$, all successes} & \multicolumn{3}{c}{$\eta_{\mathrm{inv}}$, screened} \\
\cmidrule(lr){2-4}\cmidrule(lr){5-7}
Method & median & $q_{99}$ & max & median & $q_{99}$ & max \\
\hline
\texttt{R0\_C0} & $2.69\times10^{-13}$ & $6.15\times10^{-4}$ & $6.49\times10^{-1}$
                & $3.78\times10^{-8}$ & $4.74\times10^{-2}$ & $2.45\times10^{1}$ \\
\texttt{R0\_C1} & $1.15\times10^{-15}$ & $1.76\times10^{-6}$ & $5.48\times10^{-6}$
                & $5.17\times10^{-9}$ & $9.91\times10^{-4}$ & $2.40\times10^{-3}$ \\
\texttt{R0\_C2} & $1.28\times10^{-15}$ & $2.32\times10^{-6}$ & $6.18\times10^{-6}$
                & $5.09\times10^{-9}$ & $9.91\times10^{-4}$ & $2.45\times10^{-3}$ \\
\texttt{R1\_C1} & $\mathbf{8.33\times10^{-16}}$ & $\mathbf{1.18\times10^{-6}}$ & $\mathbf{2.53\times10^{-6}}$
                & $5.09\times10^{-9}$ & $7.91\times10^{-4}$ & $\mathbf{1.48\times10^{-3}}$ \\
\texttt{R2\_C2} & $8.88\times10^{-16}$ & $1.49\times10^{-6}$ & $3.12\times10^{-6}$
                & $\mathbf{5.01\times10^{-9}}$ & $\mathbf{7.53\times10^{-4}}$ & $1.57\times10^{-3}$ \\
\hline
\end{tabular}
\end{table}

Table~\ref{tab:comparison} admits three conclusions. First, the four variants that select a pivot
by magnitude are mutually comparable and separated from the first-nonzero rule by two orders of
magnitude in the median residual and five in the maximum; the transformation defect shows the
same separation, with median $\|D_N\|_2/\|A\|_2 = 2.17\times10^{-4}$ for \texttt{R0\_C0} against
$1.18\times10^{-6}$ for \texttt{R1\_C1}. Second, the magnitude-based variants are comparable
with Gauss--Jordan inversion on the displayed median residual: all four improve on the
$1.90\times10^{-15}$ median of partial-pivoting Gauss--Jordan, while the row-pivoted variants are
also close to the $9.79\times10^{-16}$ median of complete-pivoting Gauss--Jordan. No claim of
dominance for unreported Gauss--Jordan percentiles is made. This is consistent with Fletcher's
favorable conclusion~\cite{Fletcher1997} under partial pivoting and with G\'ati's instability
finding~\cite{Gati2013} without pivoting, while emphasizing the role of the pivot strategy.
Third, among the pivoted variants there is no total ordering, the differences being within a
factor two on the statistics displayed in Table~\ref{tab:comparison}.

The large-scale campaign sharpens the third conclusion. The separation of the first-nonzero rule
persists over $14.7$ million matrices and up to $n=512$, where the median across cells of
$q_{95}(r_R/u)$ is $9.9$ and $15.4$ for \texttt{R0\_C0} in the two precisions against $1.3$ to
$2.5$ for the pivoted variants. But the ordering \emph{among} the pivoted variants is a property
of the matrix class, not of the method. Row pivoting is beneficial on random matrices with prescribed singular values, \texttt{R1\_C1}
winning the paired comparison against \texttt{R0\_C1} on $71\%$ of matrices; it is harmful on
both diagonally dominant ensembles, where \texttt{R0\_C1} wins on $85$ to $89\%$; and on
Vandermonde at $n=512$ the two magnitude rules without row pivoting are the viable choices:
\texttt{R0\_C1} has zero breakdown and \texttt{R0\_C2} breaks down on only $0.1\%$ of cases,
against $65.6\%$, $72.7\%$ and $84.6\%$ for \texttt{R0\_C0}, \texttt{R1\_C1} and
\texttt{R2\_C2}; among the two no-row-pivot rules, \texttt{R0\_C2} has the smaller conditional
$q_{95}(r_R/u)$ ($0.288$ versus $0.793$). That comparison also shows why breakdown rate and
conditional residual must be reported together:
among survivors \texttt{R2\_C2} attains $q_{95}(r_R/u)=0.296$, barely distinguishable from the
$0.288$ of \texttt{R0\_C2}, yet fails on five sixths of the sample, and where both succeed
\texttt{R0\_C2} is still the smaller on $59\%$ of matrices. The defensible statement is that
\emph{selecting a pivot by magnitude matters, and which magnitude rule is selected does not,
except through the structure of the matrix class} (SM7).

\subsection{The bound is valid but not sharp}
\label{sec:law}

The bound of Theorem~\ref{thm:backward}, evaluated with the quantities recorded during the
computation, is
%
\begin{equation}\label{eq:corrected-bound}
\Lambda_N^{\rm corr}
= \left(\frac{\|D_N\|_2}{\|A\|_2}
+ \gamma_2\sum_{j=1}^{N}\bigl\|\,|\widehat V_jP_j|\,|\widehat B_j|\,\bigr\|_2
\bigl\|\widehat{\mathcal B}_{j+1:N}\bigr\|_2\right)
\bigl\|\widehat V_{N+1}^{-1}\bigr\|_2 .
\end{equation}
%
No violation is observed over the $52\,184$ reliability-screened comparisons; the maximum observed ratio
$\eta_{\mathrm{inv}}/\Lambda_N^{\rm corr}$ being $0.266$, with $q_{50}=3.69\times10^{-6}$ and
$q_{99}=4.94\times10^{-2}$. It is, however, far too pessimistic to describe the observed
magnitude, its upper percentiles reaching $10^{52}$--$10^{61}$.

The observed loss is overwhelmingly concentrated in one place and has a structural component. The amplification factor is not the difficulty,
$\|\widehat V_{N+1}^{-1}\|_2$ having median $1.000$; nor is the defect, whose contribution
accounts for a median fraction $<10^{-4}$ of $\Lambda_N^{\rm corr}$ while tracking
$\eta_{\mathrm{inv}}$ closely. More than $99.99\%$ of the bound is the accumulation sum in the median, and
Proposition~\ref{prop:growth-floor} shows a factor of order $n^2u$ in the resulting residual bound
to be structural. The universal finite-$u$ statement is
$\Gamma_N\ge n/(1+n\gamma_2)$; in addition, the campaign happens to satisfy the stronger empirical
inequality $\Gamma_N\ge n$ on all $8320$ recorded cells, with minimum $\Gamma_N/n=1.64$.
Accordingly the theoretical right-hand side has the first-order floor $n(n+3)u+O(u^2)$, whereas
the measured $r_R/u$ has median $1.20$ over the $73\,073$ unrefined pivoted records of the
large-scale campaign. Beyond that floor the pessimism is the familiar one of bounding separately
two factors whose product telescopes. In exact arithmetic, with processed-row permutation $\Pi$,
$V_1\prod_jP_jB_j=(\Pi A)^{-1}=A^{-1}\Pi^T$; for the exactly applied computed factors their
departure from this relation is encoded exactly by $D_N=I-\Pi A\widetilde V_{N+1}$. The
individual factors carry no analogous terminal constraint. On the prescribed-singular-value
family, $\max_j\|\widehat V_j\|_2$ tracks $\|A^{-1}\|_2$ with median ratio $1.08$ and
log-correlation $0.99$, and the trailing products behave similarly. Their separate-factor product
can therefore inherit roughly squared conditioning on that family, without implying a universal
scaling law for the actual inverse backward error.

\subsection{The diagonally dominant class}
\label{sec:dd-numerics}

Theorem~\ref{thm:dd} is the one place where the analysis delivers an exact-trajectory growth
bound without a separate a priori tableau-growth hypothesis, and the large-scale campaign was
designed to test both its covered variants and the behavior of nearby variants outside its
hypotheses. Two ensembles are used. B1 has random sign patterns. For B2 a zero-diagonal row-stochastic matrix $P$ with positive off-diagonal entries is generated and
\[
A=\frac{I-\beta P}{1+\beta},
\qquad
\beta=\frac{1-\alpha}{1+\alpha}.
\]
Thus every reference row has absolute sum one and dominance gap
$(1-\beta)/(1+\beta)=\alpha$ algebraically. The campaign stores the measured reference gap as well
as the nominal grid value $\alpha\in\{10^{-3},\dots,0.5\}$; subsequent conversion to the working
precision makes the finite-precision checks numerical counterparts, not exact restatements, of the
theorem.

Across the large audit, the two theorem-covered variants \texttt{R0\_C0} and \texttt{R0\_C1}
show no computed-pivot violation: the minimum recorded value of
$\min_l|\hat t_l^{\,l}|/\alpha$ is $1.423967$, while the minimum over all five variants is
$1.419758$.
The worst recorded ratios to the nominal tableau, elementary-transformation, and exact-growth
scales are approximately $0.54$, $0.17$, and $0.038$, respectively. After synchronizing the
residual diagnostic with Proposition~\ref{prop:growth-bound}, the largest value of
\[
\frac{r_R^\infty}{\gamma_{n+3}\,6n/\alpha^3}
\]
is $1.000394\times10^{-3}$ for \texttt{R0\_C0}/\texttt{R0\_C1} and
$1.177932\times10^{-3}$ over all five variants. These finite-$u$ ratios are diagnostic comparisons;
they are not a numerical proof of the asymptotic $O(u^2)$ statement in Corollary~\ref{cor:dd-fp}.
The checks on \texttt{R0\_C2}, \texttt{R1\_C1}, and \texttt{R2\_C2} are additional empirical DD
checks outside the theorem's hypotheses, whereas those on \texttt{R0\_C0} and \texttt{R0\_C1}
are numerical counterparts of the theorem-covered paths. The audit count is deliberately not
converted into a number of independent matrix records, because several stepwise inequalities are
evaluated on each matrix.

The interchange counts delimit the exact result sharply. On B1, \texttt{R0\_C0},
\texttt{R0\_C1}, and \texttt{R0\_C2} perform zero column interchanges on every one of the
$2.6\times10^5$ balanced matrices per precision. The dedicated replay control documented in SM8
compares \texttt{R0\_C0} with \texttt{R0\_C1}: over all $520$ blocks it finds identical pivot
paths and bit-for-bit equality of $D_N$, $R_N$, and $\Pi AE_{N+1}$. No corresponding bit-for-bit
replay claim is made for \texttt{R0\_C2}. On B2 the maximum-scaled rule departs, interchanging on
$4.91\%$ of matrices, concentrated at small $n$ and small $\alpha$ and absent from $n\ge181$;
this empirical counterexample is why the theorem is not extended to \texttt{R0\_C2}. By contrast,
\texttt{R1\_C1} and \texttt{R2\_C2} permute almost every step, with median $504$ and $505$ row
interchanges at $n=512$, and those variants degrade the residual on both ensembles.

The bound is deliberately conservative in its $\alpha$ dependence. Fitting
$\Gamma_N/N\sim C\alpha^{-p}$ over the ten tested values gives effective exponents
$p\simeq0.098$ on B1 and $p\simeq0.281$ on B2 in both precisions, against the proven exponent $3$.
At $\alpha=10^{-3}$ the factor $6/\alpha^3$ exceeds the measured $\Gamma_N/N\simeq3.6$ by about
nine orders of magnitude. On B2 the local and trailing factors entering~\eqref{eq:growth} each
carry substantial $\alpha$ dependence, but most of it disappears after division by
$\|V^{N+1}\|_\infty$. This identifies the separate-factor estimates in the proof as the point a
sharper argument would have to improve; no asymptotic law in $\alpha$ is inferred from the finite
grid.

\subsection{An exactly conditioned adversarial family}
\label{sec:adversarial}

The clearest evidence that the pivot rule, not the conditioning, can govern the growth is a
family on which the conditioning is fixed at its optimum: let $n$ be even and $A$ block diagonal
with $2\times2$ blocks
%
\begin{equation}\label{eq:adversarial}
A_\varepsilon^{(2)} =
\begin{pmatrix}
\varepsilon & \sqrt{1-\varepsilon^2}\\[2pt]
\sqrt{1-\varepsilon^2} & -\varepsilon
\end{pmatrix},
\qquad 0<\varepsilon<1 .
\end{equation}
%
Each block is orthogonal and symmetric, so $A$ is orthogonal and $\kappa_2(A)=1$ exactly.
The behavior of the pivot rules on this family is not merely observed but proved.

\begin{proposition}[Exactly conditioned adversarial family]\label{prop:adversarial}
Let $n$ be even, let $0<\varepsilon<1/\sqrt2$, write $c=\sqrt{1-\varepsilon^2}$, and let $A$
be block diagonal with blocks~\eqref{eq:adversarial}. Apply~\eqref{eq:castillo-pivot} in
exact arithmetic with $V^1=I$, rows in order, and the permutation convention of
\S\ref{sec:onestep}. Then:
\begin{enumerate}
\item under the first-nonzero rule the pivot values alternate between $\varepsilon$ and
$-1/\varepsilon$, the scaled pivots between $\varepsilon$ and $1$, so that
$p^{\rm sc}_{\min}=\varepsilon$, and
\[
\max_{1\le j\le N+1}\|V^j\|_2
= \frac{\sqrt{1+c}}{\varepsilon}
= \frac{\sqrt2}{\varepsilon}\bigl(1+O(\varepsilon^2)\bigr),
\qquad\text{independently of }n;
\]
\item under any rule that selects a pivot of maximal modulus, or of maximal scaled modulus,
the scaled pivots alternate between $c$ and $1$, so that $p^{\rm sc}_{\min}=c>1/\sqrt2$, and
\[
\max_j\|V^j\|_2
=\frac{1}{\sqrt{1-\varepsilon}}
=1+\frac{\varepsilon}{2}+O(\varepsilon^2);
\]
\item in both cases $V^{N+1}=A^{-1}=A$.
\end{enumerate}
\end{proposition}

\begin{proof}
The blocks decouple.  In one block, first-nonzero selects $\varepsilon$, giving
$v_1^2=e_1/\varepsilon$ and $v_2^2=e_2-(c/\varepsilon)e_1$; the Gram matrix has squared singular
values $(1\pm c)/\varepsilon^2$.  A maximum-modulus or maximum-scaled rule instead selects $c$ and,
after the interchange,
\[
V^2=\begin{pmatrix}0&1\\1/c&-\varepsilon/c\end{pmatrix},
\qquad
(V^2)^TV^2=c^{-2}\begin{pmatrix}1&-\varepsilon\\-\varepsilon&1\end{pmatrix},
\]
whose squared singular values are $1/(1+\varepsilon)$ and $1/(1-\varepsilon)$.  The terminal
invariant gives $V^{N+1}=A^{-1}=A$.  SM10 gives the floating-point confirmation.
\end{proof}

The campaign confirms Proposition~\ref{prop:adversarial} in floating point across eight decades of
$\varepsilon$, the measured tableau growth of \texttt{R0\_C0} agreeing with
$\sqrt{1+c}/\varepsilon$ to all recorded digits. The no-row-pivoting magnitude rules agree with the exact value $(1-\varepsilon)^{-1/2}$ and therefore
tend to one as $\varepsilon\to0$; for example, \texttt{R0\_C2} gives tableau growth $1.05$
at $\varepsilon=10^{-1}$, $1.01$ at $10^{-2}$, and $1.00$ to displayed precision thereafter.
A variant with geometrically decaying block parameters reproduces the effect with repeated rather
than isolated bad pivots (SM10). Two consequences follow, and one caution. First, the tableau growth, and with
it the hypotheses of the conditional statement following~\eqref{eq:hypotheses}, is not controlled by $\kappa_2(A)$, so any
conditional statement phrased through pivot size and growth must retain a hypothesis on the pivot
sequence. Second, tableau growth is not in general a proxy for conditioning, so the close
tracking of $\max_j\|\widehat V_j\|_2$ by $\|A^{-1}\|_2$ observed in \S\ref{sec:law} must not be
read as a universal relation. The caution is that this family does not by itself decide whether
the backward error \emph{itself} admits a bound in $\kappa_2(A)$ alone, its computed residuals
being of order $\varepsilon$, so it defeats the bound and not the quantity bounded. The
repeated-bad-pivots variant adds strong finite-precision evidence for the unpivoted rule: its extreme record, at
$n=32$ in binary64 with $\varepsilon=10^{-6}$ and geometric decay $0.3$, has
%
\begin{equation}\label{eq:counterexample}
\kappa_2(A) = 1 + O(10^{-16}),
\qquad
\kappa_2(\widehat V_{N+1}) = 1,
\qquad
\eta_{\mathrm{inv}} = 6.80\times10^{-9},
\end{equation}
%
that is $\eta_{\mathrm{inv}}/u = 6.13\times10^{7}$, the largest normalized value recorded
anywhere in the campaign. At a minimum, this record shows that any $O(u)$ backward-error bound for the first-nonzero
rule whose constant is controlled only by $\kappa_2(A)$ would require a constant exceeding
$6.1\times10^7$ already at $\kappa_2(A)=1$ over the tested family. Unlike the exact tableau-growth
construction above, a finite experiment does not rule out every conceivable function of
$\kappa_2(A)$; it rules out a practically meaningful modest constant on the tested regime. The
large normalized error is not isolated: over the family the median rises from $4.1\times10^{1}$ at decay
$1.0$ to $1.6\times10^{4}$ at decay $0.1$, and from $1.4\times10^{3}$ at $\varepsilon=10^{-2}$ to
$4.5\times10^{5}$ at $\varepsilon=10^{-10}$. We claim no asymptotic law: the growth in
$\varepsilon$ is markedly sublinear, roughly $\varepsilon^{-0.3}$ over eight decades, the maxima
are not monotone, and the two parameters are not varied independently. For the pivoted variants,
whose normalized values never exceed $1.8\times10^{2}$, the question remains open.

\subsection{Empirical predictors and the explicit-inverse barrier}
\label{sec:predictors}

Over the $51\,392$ positive reliability-screened records, the spread
$\log_{10}(q_{99}/q_{01})$ of $\eta_{\rm inv}$/predictor is $6.9$ for
$\|D_N\|_2/\|A\|_2$ and $8.3$ for $un/p_{\min}^{\rm sc}$, versus $17.2$ for
$un\kappa_2(A)$ and $19.8$ for $\Lambda_N^{\rm corr}$.  Thus $D_N$ is the tightest ex-post
diagnostic tested, whereas the scaled pivot is the strongest local indicator.  The empirical law
\begin{equation}\label{eq:empirical-law}
\eta_{\rm inv}\le\frac{un}{p_{\min}^{\rm sc}}
\end{equation}
holds in $99.0\%$ of screened records but fails on $20.5\%$ of Vandermonde cases, sometimes by more
than three orders of magnitude.  Regression is compatible with the importance of the pivot, but
leave-one-family-out tests show that growth variables retain class-dependent information (SM5--SM6).

Independently of the algorithm,~\eqref{eq:deltaA} gives the explicit-inverse barrier
\begin{equation}\label{eq:barrier}
\eta_{\rm inv}\le r_R\,\kappa_2(\widehat V_{N+1}).
\end{equation}
Hence a residual of order $u$ need not imply a matrix backward error of order $u$ when the stored
inverse is ill-conditioned~\cite{DuCrozHigham1992,Higham2002}.  Our operational screen is designed
to avoid interpreting unresolved solves as measurements of $\eta_{\rm inv}$; we make no claim that
\eqref{eq:barrier} is close to equality in the campaign.

\subsection{How row pivoting degrades the residual}
\label{sec:masking}

A dedicated same-matrix replay on the two diagonally dominant ensembles recomputed $D_N$, $R_N$,
and $AE_{N+1}$ under four variants, yielding $1.04$ million complete
\texttt{R1\_C1}/\texttt{R0\_C1} pairs.  Passing from \texttt{R0\_C1} to \texttt{R1\_C1}
multiplies the geometric mean of $\|D_N\|_2$ by $2.57$ and $2.16$ on B1 and B2 in binary32, and by
$2.49$ and $2.13$ in binary64.  The corresponding ratios for $\|AE_{N+1}\|_2$ are essentially one,
and $\Gamma_N/n$ changes only marginally.  Row pivoting therefore amplifies the formation defect on
these structured classes rather than the application error.

The final residual need not inherit this amplification proportionally.  With
$H=(\|D_N\|_2^2+\|AE_{N+1}\|_2^2)^{1/2}$, $Q=\|R_N\|_2/H$, and baseline weight
$w=\|D_N\|_2^2/H^2$, the pairwise ratios satisfy exact scalar identities.  SM8 factors the residual
ratio into the defect amplification, a baseline-weight factor, an application-error factor, and a
residual-suppression factor.  The dominant multiplicative attenuation is the baseline-weight factor:
matrices on which row pivoting amplifies the defect most strongly often have small baseline defect
weight.  All ratio identities are stated on their natural positive-denominator domain; the replay
acceptance gate is absolute and operational rather than an interval certificate.

\subsection{Dimensional scaling, and the solution backward error}
\label{sec:scaling}

Two further questions are settled briefly. First, the thirteen dimensions of the large-scale
campaign allow the growth of the residual with $n$ to be estimated rather than assumed. Fitting
$q_{95}(r_R/u)\sim Cn^{s}$ per cell of family, parameters and precision, with cluster-bootstrap
intervals resampling matrices within each $n$, gives median exponents $0.22$ to $0.35$ for the
pivoted variants and $0.49$ to $0.55$ for \texttt{R0\_C0}, with median interval widths $0.011$ to
$0.014$. Two cautions are material: the variation between family--parameter cells, from $-3.6$ to $2.2$
for \texttt{R0\_C0}, is an order of magnitude larger than the statistical uncertainty, so a pooled
exponent is not meaningful; and a quadratic term in $\log n$ is preferred by AIC in $70$ to
$94\%$ of cells, with convex curvature in the majority. These are therefore \emph{effective slopes
over the recovered balanced range $8\le n\le512$}, not asymptotic exponents. Since the dimensions
above $181$ come from deterministic wall-clock truncation rather than a missing-at-random sample,
they support a finite-range trend analysis rather than a claim about the full designed population
(SM7).

Second, the normwise backward error of the computed solution,
$\eta_{\rm sol} = \|b-A\widehat x\|_2/(\|A\|_2\|\widehat x\|_2 + \|b\|_2)$, behaves differently
from $\eta_{\mathrm{inv}}$, as Proposition~\ref{prop:sol-vs-inv} predicts. Two refinement steps
bring the Castillo medians from $10^{-9}$ to $4$--$6\times10^{-15}$ while leaving
$\eta_{\mathrm{inv}}$ unchanged, the median reduction of the forward error over the large-scale
campaign being $56\%$, though several binary32 blocks show no reduction of $q_{95}$ at all. Taken
over all successful cases the $99$th percentile of $\eta_{\rm sol}$ appears to reverse the
ordering of \S\ref{sec:comparison}, ranking \texttt{R0\_C0} ahead of the pivoted variants, and
the reversal survives restriction to the configurations solved by all five; but it is confined to
the regime in which $\eta_{\rm sol}$ ceases to be informative, ninety percent of the cases with
$\eta_{\rm sol}>10^{-3}$ being binary32 with median condition number $2.3\times10^{8}$ and median
forward error $1.27$, and on the resolved subset \texttt{R0\_C0} becomes the worst variant by
factors $3.8$ and $4.4$ (SM9). The normwise solution backward error is therefore not a reliable
instrument for ranking inversion methods once the problem is unresolved in the working
precision.

\section{Consequences for the backward-error analysis}
\label{sec:discussion}

The experiments change the practical reading of the exact formulas in three ways.  First, the
transformation defect is not a small correction to a useful accumulation estimate: on the screened
records it is the sharpest ex-post diagnostic measured, whereas more than $99.99\%$ of the full
upper bound comes from the accumulation term in the median.  A sharper theory should therefore
control the formation of the elementary transformations without replacing the trailing products by
separate worst-case norms at every step.

Second, stability is conditional and the local bottleneck is the pivot together with dot-product
sensitivity.  The empirical pivot relation succeeds on most tested records but fails systematically
on Vandermonde matrices, where cancellation matters.  The exact orthogonal family proves that no
bound on intermediate tableau growth can depend on $\kappa_2(A)$ alone; the repeated-bad-pivot
family adds a finite-precision warning for the backward error itself, without proving the
nonexistence of every possible condition-number-only function.  Conversely, strict row diagonal
dominance protects the active pivot and forces first-nonzero and maximum-modulus onto the same benign
path.  Magnitude-based selection is therefore an essential safeguard on the unstructured and
adversarial classes tested here, but no single magnitude rule is uniformly best.

Third, inverse stability and solution stability must be separated.  Equation~\eqref{eq:sol-residual}
shows that a right-hand side probes one direction of $R_N$ and adds only the rounding of one
matrix--vector product; iterative refinement changes this solution error without changing the stored
inverse.  The apparent ranking reversal in $\eta_{\rm sol}$ occurs only in the unresolved regime and
disappears once $\kappa_2(A)u<10^{-2}$ is imposed (SM9).

\section{Conclusions}
\label{sec:conclusions}

The explicit Purcell--Egerv\'ary--Castillo inverse has two distinct rounding mechanisms.  The exact
identity $R_N=D_N-AE_{N+1}$ separates the transformation defect created while the factors are formed
from the application error accumulated while they are applied, and
$\Delta A=R_N\widehat V_{N+1}^{-1}$ converts the residual into the inverse backward error.  The
resulting growth factor has a universal linear floor, while strict row diagonal dominance yields the
exact bound $\Gamma_N\le6n/\alpha^3$ and a stable first-nonzero/maximum-modulus pivot path.

Numerically, the classical accumulation estimate is rigorous but far too pessimistic; the defect is
the sharpest ex-post diagnostic measured and the scaled pivot is the dominant local indicator.  The
orthogonal construction proves that intermediate growth cannot be controlled by $\kappa_2(A)$
alone, while the campaigns show that magnitude-based pivoting is an important safeguard on the
unstructured and adversarial classes tested.  No unconditional backward-stability theorem follows,
and no magnitude rule is uniformly best.  A sharper theory should propagate the local formation
defect while preserving cancellations in the trailing products and retaining both pivot size and
dot-product sensitivity.

\par\medskip
\noindent\textbf{Data and code availability.} Code and reproducibility data are available at \href{https://doi.org/10.5281/zenodo.21794454}{doi:10.5281/zenodo.21794454}.
\par

\par\medskip
\noindent\textbf{Acknowledgments.}
The author acknowledges the technical expertise and assistance provided by the SCBI (Supercomputing and Bioinformatics Center) at the Universidad de M\'alaga, including access to the Picasso Supercomputer.
\par

\bibliographystyle{siamplain}
\bibliography{castillo_refs}

\end{document}
